% Options for packages loaded elsewhere
\PassOptionsToPackage{unicode}{hyperref}
\PassOptionsToPackage{hyphens}{url}
\documentclass[
]{article}
\usepackage{xcolor}
\usepackage[margin=1in]{geometry}
\usepackage{amsmath,amssymb}
\setcounter{secnumdepth}{-\maxdimen} % remove section numbering
\usepackage{iftex}
\ifPDFTeX
  \usepackage[T1]{fontenc}
  \usepackage[utf8]{inputenc}
  \usepackage{textcomp} % provide euro and other symbols
\else % if luatex or xetex
  \usepackage{unicode-math} % this also loads fontspec
  \defaultfontfeatures{Scale=MatchLowercase}
  \defaultfontfeatures[\rmfamily]{Ligatures=TeX,Scale=1}
\fi
\usepackage{lmodern}
\ifPDFTeX\else
  % xetex/luatex font selection
\fi
% Use upquote if available, for straight quotes in verbatim environments
\IfFileExists{upquote.sty}{\usepackage{upquote}}{}
\IfFileExists{microtype.sty}{% use microtype if available
  \usepackage[]{microtype}
  \UseMicrotypeSet[protrusion]{basicmath} % disable protrusion for tt fonts
}{}
\makeatletter
\@ifundefined{KOMAClassName}{% if non-KOMA class
  \IfFileExists{parskip.sty}{%
    \usepackage{parskip}
  }{% else
    \setlength{\parindent}{0pt}
    \setlength{\parskip}{6pt plus 2pt minus 1pt}}
}{% if KOMA class
  \KOMAoptions{parskip=half}}
\makeatother
\usepackage{color}
\usepackage{fancyvrb}
\newcommand{\VerbBar}{|}
\newcommand{\VERB}{\Verb[commandchars=\\\{\}]}
\DefineVerbatimEnvironment{Highlighting}{Verbatim}{commandchars=\\\{\}}
% Add ',fontsize=\small' for more characters per line
\usepackage{framed}
\definecolor{shadecolor}{RGB}{248,248,248}
\newenvironment{Shaded}{\begin{snugshade}}{\end{snugshade}}
\newcommand{\AlertTok}[1]{\textcolor[rgb]{0.94,0.16,0.16}{#1}}
\newcommand{\AnnotationTok}[1]{\textcolor[rgb]{0.56,0.35,0.01}{\textbf{\textit{#1}}}}
\newcommand{\AttributeTok}[1]{\textcolor[rgb]{0.13,0.29,0.53}{#1}}
\newcommand{\BaseNTok}[1]{\textcolor[rgb]{0.00,0.00,0.81}{#1}}
\newcommand{\BuiltInTok}[1]{#1}
\newcommand{\CharTok}[1]{\textcolor[rgb]{0.31,0.60,0.02}{#1}}
\newcommand{\CommentTok}[1]{\textcolor[rgb]{0.56,0.35,0.01}{\textit{#1}}}
\newcommand{\CommentVarTok}[1]{\textcolor[rgb]{0.56,0.35,0.01}{\textbf{\textit{#1}}}}
\newcommand{\ConstantTok}[1]{\textcolor[rgb]{0.56,0.35,0.01}{#1}}
\newcommand{\ControlFlowTok}[1]{\textcolor[rgb]{0.13,0.29,0.53}{\textbf{#1}}}
\newcommand{\DataTypeTok}[1]{\textcolor[rgb]{0.13,0.29,0.53}{#1}}
\newcommand{\DecValTok}[1]{\textcolor[rgb]{0.00,0.00,0.81}{#1}}
\newcommand{\DocumentationTok}[1]{\textcolor[rgb]{0.56,0.35,0.01}{\textbf{\textit{#1}}}}
\newcommand{\ErrorTok}[1]{\textcolor[rgb]{0.64,0.00,0.00}{\textbf{#1}}}
\newcommand{\ExtensionTok}[1]{#1}
\newcommand{\FloatTok}[1]{\textcolor[rgb]{0.00,0.00,0.81}{#1}}
\newcommand{\FunctionTok}[1]{\textcolor[rgb]{0.13,0.29,0.53}{\textbf{#1}}}
\newcommand{\ImportTok}[1]{#1}
\newcommand{\InformationTok}[1]{\textcolor[rgb]{0.56,0.35,0.01}{\textbf{\textit{#1}}}}
\newcommand{\KeywordTok}[1]{\textcolor[rgb]{0.13,0.29,0.53}{\textbf{#1}}}
\newcommand{\NormalTok}[1]{#1}
\newcommand{\OperatorTok}[1]{\textcolor[rgb]{0.81,0.36,0.00}{\textbf{#1}}}
\newcommand{\OtherTok}[1]{\textcolor[rgb]{0.56,0.35,0.01}{#1}}
\newcommand{\PreprocessorTok}[1]{\textcolor[rgb]{0.56,0.35,0.01}{\textit{#1}}}
\newcommand{\RegionMarkerTok}[1]{#1}
\newcommand{\SpecialCharTok}[1]{\textcolor[rgb]{0.81,0.36,0.00}{\textbf{#1}}}
\newcommand{\SpecialStringTok}[1]{\textcolor[rgb]{0.31,0.60,0.02}{#1}}
\newcommand{\StringTok}[1]{\textcolor[rgb]{0.31,0.60,0.02}{#1}}
\newcommand{\VariableTok}[1]{\textcolor[rgb]{0.00,0.00,0.00}{#1}}
\newcommand{\VerbatimStringTok}[1]{\textcolor[rgb]{0.31,0.60,0.02}{#1}}
\newcommand{\WarningTok}[1]{\textcolor[rgb]{0.56,0.35,0.01}{\textbf{\textit{#1}}}}
\usepackage{graphicx}
\makeatletter
\newsavebox\pandoc@box
\newcommand*\pandocbounded[1]{% scales image to fit in text height/width
  \sbox\pandoc@box{#1}%
  \Gscale@div\@tempa{\textheight}{\dimexpr\ht\pandoc@box+\dp\pandoc@box\relax}%
  \Gscale@div\@tempb{\linewidth}{\wd\pandoc@box}%
  \ifdim\@tempb\p@<\@tempa\p@\let\@tempa\@tempb\fi% select the smaller of both
  \ifdim\@tempa\p@<\p@\scalebox{\@tempa}{\usebox\pandoc@box}%
  \else\usebox{\pandoc@box}%
  \fi%
}
% Set default figure placement to htbp
\def\fps@figure{htbp}
\makeatother
\setlength{\emergencystretch}{3em} % prevent overfull lines
\providecommand{\tightlist}{%
  \setlength{\itemsep}{0pt}\setlength{\parskip}{0pt}}
\usepackage{bookmark}
\IfFileExists{xurl.sty}{\usepackage{xurl}}{} % add URL line breaks if available
\urlstyle{same}
\hypersetup{
  hidelinks,
  pdfcreator={LaTeX via pandoc}}

\author{}
\date{\vspace{-2.5em}}

\begin{document}

\begin{center}\rule{0.5\linewidth}{0.5pt}\end{center}

\begin{center}\rule{0.5\linewidth}{0.5pt}\end{center}

title: ``The relationship between GDP per Capita and Unemployment''

author: ``Tutorial 5 group 3''

date: ``2026-06-22''

output: pdf\_document

\begin{center}\rule{0.5\linewidth}{0.5pt}\end{center}

\section{Set-up your environment}\label{set-up-your-environment}

\begin{Shaded}
\begin{Highlighting}[]
\FunctionTok{require}\NormalTok{(tidyverse)}
\end{Highlighting}
\end{Shaded}

\begin{verbatim}
## Loading required package: tidyverse
\end{verbatim}

\begin{verbatim}
## -- Attaching core tidyverse packages ---------------------------------- tidyverse 2.0.0 --
## v dplyr     1.2.1     v readr     2.2.0
## v forcats   1.0.1     v stringr   1.6.0
## v ggplot2   4.0.3     v tibble    3.3.1
## v lubridate 1.9.5     v tidyr     1.3.2
## v purrr     1.2.2     
## -- Conflicts ---------------------------------------------------- tidyverse_conflicts() --
## x dplyr::filter() masks stats::filter()
## x dplyr::lag()    masks stats::lag()
## i Use the conflicted package (<http://conflicted.r-lib.org/>) to force all conflicts to become errors
\end{verbatim}

\begin{Shaded}
\begin{Highlighting}[]
\FunctionTok{require}\NormalTok{(rmarkdown)}
\end{Highlighting}
\end{Shaded}

\begin{verbatim}
## Loading required package: rmarkdown
\end{verbatim}

\begin{Shaded}
\begin{Highlighting}[]
\FunctionTok{require}\NormalTok{(yaml)}
\end{Highlighting}
\end{Shaded}

\begin{verbatim}
## Loading required package: yaml
\end{verbatim}

\begin{Shaded}
\begin{Highlighting}[]
\FunctionTok{require}\NormalTok{(readxl)}
\end{Highlighting}
\end{Shaded}

\begin{verbatim}
## Loading required package: readxl
\end{verbatim}

\begin{Shaded}
\begin{Highlighting}[]
\FunctionTok{library}\NormalTok{(dplyr)}
\FunctionTok{library}\NormalTok{(ggplot2)}
\end{Highlighting}
\end{Shaded}

\section{Title Page}\label{title-page}

Teammembers:\\
Koen Damen (2889378)

Thomas Dop (2898269)

Tim de Marez Oyens (2882348)

Timo Bakker (2863621)

Daniel Ferrier (2900731)

Clémence van Gils (2898031)

Tutorial 5 Group 3

Tutorial lecturer: Simon Loop

\section{Part 1 - Identify a Social
Problem}\label{part-1---identify-a-social-problem}

\subsection{1.1 The Relationship Between GDP Growth and
Unemployment}\label{the-relationship-between-gdp-growth-and-unemployment}

GDP is one of the most important economic indicators. It measures the
total value of all goods and services produced within a country.
Governments often use GDP as a measure of how well a society is doing.

According to Okun's law, there is a negative relationship between
economic growth and unemployment: if economic activity falls, there is
less production and labor demand causing people to lose their jobs. This
will lead to a rise in the unemployment rate. Conversely, if the economy
expands and GDP grows, unemployment will fall, as investment increases,
consumption rises and, as a result, more people are hired. (Blanchard et
al., 2021) Furthermore, high unemployment will bring about high costs
for the government.

That is why it is important to conduct research to gain insight into the
relationship between unemployment and GDP. By examining this
relationship, we can determine the extent to which economic growth
contributes to increased employment. The findings could help
policymakers to take measures that stimulate economic growth and also
help to reduce unemployment. In this report we quantify the the
relationship between the GDP per capita (\emph{Gross Domestic Product
(GDP) per Capita - Data Portal - United Nations Economic Commission for
Europe}, n.d.) and the unemployment rate (\emph{Unemployment Rate - Data
Portal - United Nations Economic Commission for Europe}, n.d.) from 2005
until 2024.\\

\section{Part 2 - Data Sourcing}\label{part-2---data-sourcing}

\subsection{2.1 Load in the data}\label{load-in-the-data}

\subsection{2.2 Provide a short summary of the
dataset(s)}\label{provide-a-short-summary-of-the-datasets}

Dataset 1 contains the procentual unemployment rate for 55 countries
from 2005 to 2025. Dataset 2 contains the GDP per capita in USD adjusted
by purchasing power parities. We have data for 56 countries for the
years 2005 to 2024. For some countries we are missing some years of GDP
and unemployment value.

\begin{Shaded}
\begin{Highlighting}[]
\FunctionTok{names}\NormalTok{(dataset1)}
\end{Highlighting}
\end{Shaded}

\begin{verbatim}
##  [1] "X.1"             "IndicatorCode"   "IndicatorName"   "VariableName"   
##  [5] "MeasurementName" "CountryCode"     "Alpha3Code"      "CountryName"    
##  [9] "PeriodCode"      "Value"           "X"
\end{verbatim}

\begin{Shaded}
\begin{Highlighting}[]
\FunctionTok{names}\NormalTok{(dataset2)}
\end{Highlighting}
\end{Shaded}

\begin{verbatim}
##  [1] "X.1"             "IndicatorCode"   "IndicatorName"   "VariableName"   
##  [5] "MeasurementName" "CountryCode"     "Alpha3Code"      "CountryName"    
##  [9] "PeriodCode"      "Value"           "X"
\end{verbatim}

\begin{Shaded}
\begin{Highlighting}[]
\FunctionTok{head}\NormalTok{(dataset1)}
\end{Highlighting}
\end{Shaded}

\begin{verbatim}
##   X.1 IndicatorCode        IndicatorName      VariableName MeasurementName CountryCode
## 1   1            43 Unemployment rate, % Unemployment rate         Percent           8
## 2   2            43 Unemployment rate, % Unemployment rate         Percent           8
## 3   3            43 Unemployment rate, % Unemployment rate         Percent           8
## 4   4            43 Unemployment rate, % Unemployment rate         Percent           8
## 5   5            43 Unemployment rate, % Unemployment rate         Percent           8
## 6   6            43 Unemployment rate, % Unemployment rate         Percent           8
##   Alpha3Code CountryName PeriodCode Value  X
## 1        ALB     Albania       2005  14.3 NA
## 2        ALB     Albania       2006  13.9 NA
## 3        ALB     Albania       2007  13.4 NA
## 4        ALB     Albania       2008  12.8 NA
## 5        ALB     Albania       2009  13.6 NA
## 6        ALB     Albania       2010  13.7 NA
\end{verbatim}

\begin{Shaded}
\begin{Highlighting}[]
\FunctionTok{head}\NormalTok{(dataset2)}
\end{Highlighting}
\end{Shaded}

\begin{verbatim}
##   X.1 IndicatorCode                                               IndicatorName
## 1   1            12 GDP per capita in USD adjusted by purchasing power parities
## 2   2            12 GDP per capita in USD adjusted by purchasing power parities
## 3   3            12 GDP per capita in USD adjusted by purchasing power parities
## 4   4            12 GDP per capita in USD adjusted by purchasing power parities
## 5   5            12 GDP per capita in USD adjusted by purchasing power parities
## 6   6            12 GDP per capita in USD adjusted by purchasing power parities
##                              VariableName                                 MeasurementName
## 1 Gross domestic product (GDP) per capita Total USD adjusted by purchasing power parities
## 2 Gross domestic product (GDP) per capita Total USD adjusted by purchasing power parities
## 3 Gross domestic product (GDP) per capita Total USD adjusted by purchasing power parities
## 4 Gross domestic product (GDP) per capita Total USD adjusted by purchasing power parities
## 5 Gross domestic product (GDP) per capita Total USD adjusted by purchasing power parities
## 6 Gross domestic product (GDP) per capita Total USD adjusted by purchasing power parities
##   CountryCode Alpha3Code CountryName PeriodCode  Value  X
## 1           8        ALB     Albania       2005 6014.3 NA
## 2           8        ALB     Albania       2006 6754.3 NA
## 3           8        ALB     Albania       2007 7585.0 NA
## 4           8        ALB     Albania       2008 8469.3 NA
## 5           8        ALB     Albania       2009 9025.9 NA
## 6           8        ALB     Albania       2010 9755.9 NA
\end{verbatim}

\begin{Shaded}
\begin{Highlighting}[]
\FunctionTok{dim}\NormalTok{(dataset1)}
\end{Highlighting}
\end{Shaded}

\begin{verbatim}
## [1] 1155   11
\end{verbatim}

\begin{Shaded}
\begin{Highlighting}[]
\FunctionTok{dim}\NormalTok{(dataset2)}
\end{Highlighting}
\end{Shaded}

\begin{verbatim}
## [1] 1120   11
\end{verbatim}

\begin{Shaded}
\begin{Highlighting}[]
\NormalTok{inline\_code }\OtherTok{=} \ConstantTok{TRUE}
\end{Highlighting}
\end{Shaded}

\subsection{2.3 Describe the type of variables
included}\label{describe-the-type-of-variables-included}

The variables contain information about the socioeconomic status (SES)
from citizens drawn from administrative and statistical sources from the
UNECE database. GDP per capita reflects economic performance and
unemployment rate reflects the situation in the labour market.

\textbf{Dataset 1}~

Dataset 1 is a panel dataset with measurements from 55 countries from
2005 to 2025.\\
The main variable is the unemployment rate, measured as a percentage of
the total labour force. This is a continuous quantitative variable.\\
Countries are identified by categorical variables such as CountryCode
and Alpha3Code. The variable PeriodCode is the identifier for the time
dimension, shown in years.\\
The metadata variables, like Indicator Code and Variable name, are used
to show what is being measured by set variables.

\textbf{Dataset 2}~\\
Dataset 2 is also a panel dataset, containing measurements from 56
countries from 2005 to 2024.\\
The main variable is GDP per capita, measured in PPP-adjusted US
dollars. This is a continuous quantitative variable.\\
Similar to dataset1, countries are identified by categorical variables
and the PeriodCode tracks the time dimension.\\
The metadata variables, like Indicator Code and Variable name, are used
to show what is being measured by set variables.

\section{Part 3 - Quantifying}\label{part-3---quantifying}

\subsection{3.1 Data cleaning}\label{data-cleaning}

\begin{Shaded}
\begin{Highlighting}[]
\NormalTok{dataset1\_cleaned }\OtherTok{\textless{}{-}}\NormalTok{ dataset1 }\SpecialCharTok{\%\textgreater{}\%}
  \FunctionTok{select}\NormalTok{(IndicatorName, CountryName, PeriodCode , Value) }


\NormalTok{dataset2\_cleaned }\OtherTok{\textless{}{-}}\NormalTok{ dataset2 }\SpecialCharTok{\%\textgreater{}\%} 
  \FunctionTok{select}\NormalTok{(IndicatorName, CountryName, PeriodCode, Value)}

\NormalTok{merged\_data }\OtherTok{\textless{}{-}} \FunctionTok{inner\_join}\NormalTok{(}
\NormalTok{  dataset1\_cleaned,}
\NormalTok{  dataset2\_cleaned,}
  \AttributeTok{by =} \FunctionTok{c}\NormalTok{(}\StringTok{"CountryName"}\NormalTok{, }\StringTok{"PeriodCode"}\NormalTok{)}
\NormalTok{)}
\end{Highlighting}
\end{Shaded}

\subsection{3.2}\label{section}

Variable 1

\begin{Shaded}
\begin{Highlighting}[]
\FunctionTok{library}\NormalTok{(dplyr)}

\NormalTok{lagged\_gdp\_table }\OtherTok{\textless{}{-}}\NormalTok{ merged\_data }\SpecialCharTok{\%\textgreater{}\%}
  \FunctionTok{arrange}\NormalTok{(CountryName, PeriodCode) }\SpecialCharTok{\%\textgreater{}\%}
  \FunctionTok{group\_by}\NormalTok{(CountryName) }\SpecialCharTok{\%\textgreater{}\%}
  \FunctionTok{mutate}\NormalTok{(}
    \CommentTok{\# Force Value.x to be numeric so the math works}
    \AttributeTok{Value.x =} \FunctionTok{as.numeric}\NormalTok{(Value.x),}
    
    \CommentTok{\# Calculate GDP growth and its lag}
    \AttributeTok{GDPGrowth =}\NormalTok{ ((Value.x }\SpecialCharTok{{-}} \FunctionTok{lag}\NormalTok{(Value.x)) }\SpecialCharTok{/} \FunctionTok{lag}\NormalTok{(Value.x)) }\SpecialCharTok{*} \DecValTok{100}\NormalTok{,}
    \AttributeTok{LaggedGDPGrowth =} \FunctionTok{lag}\NormalTok{(GDPGrowth)}
\NormalTok{  ) }\SpecialCharTok{\%\textgreater{}\%}
  \FunctionTok{select}\NormalTok{(}
\NormalTok{    CountryName,}
\NormalTok{    PeriodCode,}
\NormalTok{    Value.x,}
\NormalTok{    GDPGrowth,}
\NormalTok{    LaggedGDPGrowth}
\NormalTok{  ) }\SpecialCharTok{\%\textgreater{}\%}
  \FunctionTok{ungroup}\NormalTok{()}
\end{Highlighting}
\end{Shaded}

\begin{verbatim}
## Warning: There were 10 warnings in `mutate()`.
## The first warning was:
## i In argument: `Value.x = as.numeric(Value.x)`.
## i In group 2: `CountryName = "Andorra"`.
## Caused by warning:
## ! NAs introduced by coercion
## i Run `dplyr::last_dplyr_warnings()` to see the 9 remaining warnings.
\end{verbatim}

\begin{Shaded}
\begin{Highlighting}[]
\FunctionTok{head}\NormalTok{(lagged\_gdp\_table)}
\end{Highlighting}
\end{Shaded}

\begin{verbatim}
## # A tibble: 6 x 5
##   CountryName PeriodCode Value.x GDPGrowth LaggedGDPGrowth
##   <chr>            <int>   <dbl>     <dbl>           <dbl>
## 1 Albania           2005    14.3    NA               NA   
## 2 Albania           2006    13.9    -2.80            NA   
## 3 Albania           2007    13.4    -3.60            -2.80
## 4 Albania           2008    12.8    -4.48            -3.60
## 5 Albania           2009    13.6     6.25            -4.48
## 6 Albania           2010    13.7     0.735            6.25
\end{verbatim}

\begin{verbatim}
\end{verbatim}

Variable 2

\begin{Shaded}
\begin{Highlighting}[]
\FunctionTok{library}\NormalTok{(dplyr)}

\NormalTok{elasticity\_table }\OtherTok{\textless{}{-}}\NormalTok{ merged\_data }\SpecialCharTok{\%\textgreater{}\%}
  \FunctionTok{arrange}\NormalTok{(CountryName, PeriodCode) }\SpecialCharTok{\%\textgreater{}\%}
  \FunctionTok{group\_by}\NormalTok{(CountryName) }\SpecialCharTok{\%\textgreater{}\%}
  \FunctionTok{mutate}\NormalTok{(}
    \CommentTok{\# Ensure both columns are treated as numbers}
    \AttributeTok{Value.x =} \FunctionTok{as.numeric}\NormalTok{(Value.x),}
    \AttributeTok{Value.y =} \FunctionTok{as.numeric}\NormalTok{(Value.y),}
    
    \CommentTok{\# Now run your calculations safely}
    \AttributeTok{GDPGrowth =}\NormalTok{ ((Value.y }\SpecialCharTok{{-}} \FunctionTok{lag}\NormalTok{(Value.y)) }\SpecialCharTok{/} \FunctionTok{lag}\NormalTok{(Value.y)) }\SpecialCharTok{*} \DecValTok{100}\NormalTok{,}
    \AttributeTok{UnemploymentChange =}\NormalTok{ Value.x }\SpecialCharTok{{-}} \FunctionTok{lag}\NormalTok{(Value.x),}
    \AttributeTok{Elasticity =}\NormalTok{ UnemploymentChange }\SpecialCharTok{/}\NormalTok{ GDPGrowth}
\NormalTok{  ) }\SpecialCharTok{\%\textgreater{}\%}
  \FunctionTok{ungroup}\NormalTok{()}
\end{Highlighting}
\end{Shaded}

\begin{verbatim}
## Warning: There were 21 warnings in `mutate()`.
## The first warning was:
## i In argument: `Value.x = as.numeric(Value.x)`.
## i In group 2: `CountryName = "Andorra"`.
## Caused by warning:
## ! NAs introduced by coercion
## i Run `dplyr::last_dplyr_warnings()` to see the 20 remaining warnings.
\end{verbatim}

\begin{Shaded}
\begin{Highlighting}[]
\FunctionTok{head}\NormalTok{(elasticity\_table)}
\end{Highlighting}
\end{Shaded}

\begin{verbatim}
## # A tibble: 6 x 9
##   IndicatorName.x      CountryName PeriodCode Value.x IndicatorName.y    Value.y GDPGrowth
##   <chr>                <chr>            <int>   <dbl> <chr>                <dbl>     <dbl>
## 1 Unemployment rate, % Albania           2005    14.3 GDP per capita in~   6014.     NA   
## 2 Unemployment rate, % Albania           2006    13.9 GDP per capita in~   6754.     12.3 
## 3 Unemployment rate, % Albania           2007    13.4 GDP per capita in~   7585      12.3 
## 4 Unemployment rate, % Albania           2008    12.8 GDP per capita in~   8469.     11.7 
## 5 Unemployment rate, % Albania           2009    13.6 GDP per capita in~   9026.      6.57
## 6 Unemployment rate, % Albania           2010    13.7 GDP per capita in~   9756.      8.09
## # i 2 more variables: UnemploymentChange <dbl>, Elasticity <dbl>
\end{verbatim}

\subsection{3.3 Visualize temporal
variation}\label{visualize-temporal-variation}

\begin{Shaded}
\begin{Highlighting}[]
\FunctionTok{library}\NormalTok{(ggplot2)}
\FunctionTok{library}\NormalTok{(dplyr)}

\NormalTok{plot\_data }\OtherTok{\textless{}{-}}\NormalTok{ elasticity\_table }\SpecialCharTok{\%\textgreater{}\%}
  \FunctionTok{group\_by}\NormalTok{(PeriodCode) }\SpecialCharTok{\%\textgreater{}\%}
  \FunctionTok{summarise}\NormalTok{(}
    \CommentTok{\# CHANGE \textquotesingle{}Value.y\textquotesingle{} to whatever your actual GDP column is named}
    \StringTok{\textasciigrave{}}\AttributeTok{Avg GDP per Capita}\StringTok{\textasciigrave{}} \OtherTok{=} \FunctionTok{mean}\NormalTok{(Value.y, }\AttributeTok{na.rm =} \ConstantTok{TRUE}\NormalTok{), }
    
    \CommentTok{\# CHANGE \textquotesingle{}Value.x\textquotesingle{} to whatever your actual Unemployment column is named}
    \StringTok{\textasciigrave{}}\AttributeTok{Avg Unemployment Rate}\StringTok{\textasciigrave{}} \OtherTok{=} \FunctionTok{mean}\NormalTok{(Value.x, }\AttributeTok{na.rm =} \ConstantTok{TRUE}\NormalTok{)}
\NormalTok{  ) }\SpecialCharTok{\%\textgreater{}\%}
\NormalTok{  tidyr}\SpecialCharTok{::}\FunctionTok{pivot\_longer}\NormalTok{(}
    \AttributeTok{cols =} \FunctionTok{c}\NormalTok{(}\StringTok{\textasciigrave{}}\AttributeTok{Avg GDP per Capita}\StringTok{\textasciigrave{}}\NormalTok{, }\StringTok{\textasciigrave{}}\AttributeTok{Avg Unemployment Rate}\StringTok{\textasciigrave{}}\NormalTok{), }
    \AttributeTok{names\_to =} \StringTok{"Metric"}\NormalTok{, }
    \AttributeTok{values\_to =} \StringTok{"Value"}
\NormalTok{  )}

\CommentTok{\# 2. Plot}
\FunctionTok{ggplot}\NormalTok{(plot\_data, }\FunctionTok{aes}\NormalTok{(}\AttributeTok{x =}\NormalTok{ PeriodCode, }\AttributeTok{y =}\NormalTok{ Value, }\AttributeTok{color =}\NormalTok{ Metric)) }\SpecialCharTok{+}
  \FunctionTok{geom\_line}\NormalTok{(}\AttributeTok{size =} \DecValTok{1}\NormalTok{) }\SpecialCharTok{+}
  \FunctionTok{geom\_point}\NormalTok{() }\SpecialCharTok{+}
  \FunctionTok{facet\_wrap}\NormalTok{(}\SpecialCharTok{\textasciitilde{}}\NormalTok{Metric, }\AttributeTok{ncol =} \DecValTok{1}\NormalTok{, }\AttributeTok{scales =} \StringTok{"free\_y"}\NormalTok{) }\SpecialCharTok{+} 
  \FunctionTok{scale\_x\_continuous}\NormalTok{(}\AttributeTok{breaks =} \FunctionTok{seq}\NormalTok{(}\DecValTok{2004}\NormalTok{, }\DecValTok{2024}\NormalTok{, }\AttributeTok{by =} \DecValTok{2}\NormalTok{)) }\SpecialCharTok{+} 
  \FunctionTok{labs}\NormalTok{(}
    \AttributeTok{title =} \StringTok{"Temporal Variation of GDP per Capita and Unemployment"}\NormalTok{,}
    \AttributeTok{x =} \StringTok{"Year"}\NormalTok{,}
    \AttributeTok{y =} \StringTok{"Average Value"}
\NormalTok{  ) }\SpecialCharTok{+}
  \FunctionTok{theme\_minimal}\NormalTok{() }\SpecialCharTok{+}
  \FunctionTok{theme}\NormalTok{(}\AttributeTok{legend.position =} \StringTok{"none"}\NormalTok{)}
\end{Highlighting}
\end{Shaded}

\begin{verbatim}
## Warning: Using `size` aesthetic for lines was deprecated in ggplot2 3.4.0.
## i Please use `linewidth` instead.
## This warning is displayed once per session.
## Call `lifecycle::last_lifecycle_warnings()` to see where this warning was generated.
\end{verbatim}

\pandocbounded{\includegraphics[keepaspectratio]{Template_Assignment_files/Template_Assignment_files/figure-latex/unnamed-chunk-4-1.pdf}}

\subsection{3.4 Visualize spatial
variation}\label{visualize-spatial-variation}

\begin{Shaded}
\begin{Highlighting}[]
\FunctionTok{library}\NormalTok{(dplyr)}
\FunctionTok{library}\NormalTok{(ggplot2)}

\NormalTok{merged\_data}\SpecialCharTok{$}\NormalTok{Value.x }\OtherTok{\textless{}{-}} \FunctionTok{as.numeric}\NormalTok{(merged\_data}\SpecialCharTok{$}\NormalTok{Value.x)}
\end{Highlighting}
\end{Shaded}

\begin{verbatim}
## Warning: NAs introduced by coercion
\end{verbatim}

\begin{Shaded}
\begin{Highlighting}[]
\NormalTok{spatial\_data }\OtherTok{\textless{}{-}}\NormalTok{ merged\_data }\SpecialCharTok{\%\textgreater{}\%}
  \FunctionTok{group\_by}\NormalTok{(CountryName) }\SpecialCharTok{\%\textgreater{}\%}
  \FunctionTok{summarise}\NormalTok{(}
    \AttributeTok{Avg\_Unemployment =} \FunctionTok{mean}\NormalTok{(Value.x, }\AttributeTok{na.rm =} \ConstantTok{TRUE}\NormalTok{)}
\NormalTok{  ) }\SpecialCharTok{\%\textgreater{}\%}
  \FunctionTok{arrange}\NormalTok{(}\FunctionTok{desc}\NormalTok{(Avg\_Unemployment)) }\SpecialCharTok{\%\textgreater{}\%}
  \FunctionTok{slice}\NormalTok{(}\DecValTok{1}\SpecialCharTok{:}\DecValTok{10}\NormalTok{)}

\FunctionTok{ggplot}\NormalTok{(}
\NormalTok{  spatial\_data,}
  \FunctionTok{aes}\NormalTok{(}
    \AttributeTok{x =} \FunctionTok{reorder}\NormalTok{(CountryName, Avg\_Unemployment),}
    \AttributeTok{y =}\NormalTok{ Avg\_Unemployment}
\NormalTok{  )}
\NormalTok{) }\SpecialCharTok{+}
  \FunctionTok{geom\_col}\NormalTok{() }\SpecialCharTok{+}
  \FunctionTok{coord\_flip}\NormalTok{() }\SpecialCharTok{+}
  \FunctionTok{labs}\NormalTok{(}
    \AttributeTok{title =} \StringTok{"Countries with the Highest Average Unemployment Rate"}\NormalTok{,}
    \AttributeTok{x =} \StringTok{"Country"}\NormalTok{,}
    \AttributeTok{y =} \StringTok{"Average Unemployment Rate (\%)"}
\NormalTok{  ) }\SpecialCharTok{+}
  \FunctionTok{theme\_minimal}\NormalTok{()}
\end{Highlighting}
\end{Shaded}

\pandocbounded{\includegraphics[keepaspectratio]{Template_Assignment_files/Template_Assignment_files/figure-latex/visualise_map-1.pdf}}

This figure shows the differences in unemployment rates between
countries. We can see that unemployment is not the same everywhere.
Countries such as North Macedonia, Bosnia and Herzegovina, and Georgia
have much higher unemployment rates than other countries in the dataset.
This is important for our social problem because high unemployment can
lead to poverty, financial problems, and a lower quality of life. The
graph shows that some countries are affected by unemployment much more
than others.

\subsection{3.5 Visualize sub-population
variation}\label{visualize-sub-population-variation}

\begin{Shaded}
\begin{Highlighting}[]
\NormalTok{boxplot\_data }\OtherTok{\textless{}{-}}\NormalTok{ elasticity\_table }\SpecialCharTok{\%\textgreater{}\%}
  \FunctionTok{filter}\NormalTok{(}\SpecialCharTok{!}\FunctionTok{is.na}\NormalTok{(Value.y), }\SpecialCharTok{!}\FunctionTok{is.na}\NormalTok{(Value.x)) }\SpecialCharTok{\%\textgreater{}\%}
  \FunctionTok{mutate}\NormalTok{(}
    \AttributeTok{GDP\_Group =} \FunctionTok{ifelse}\NormalTok{(}
\NormalTok{      Value.y }\SpecialCharTok{\textgreater{}=} \FunctionTok{median}\NormalTok{(Value.y, }\AttributeTok{na.rm =} \ConstantTok{TRUE}\NormalTok{),}
      \StringTok{"High GDP countries"}\NormalTok{,}
      \StringTok{"Low GDP countries"}
\NormalTok{    )}
\NormalTok{  )}

\FunctionTok{ggplot}\NormalTok{(boxplot\_data, }\FunctionTok{aes}\NormalTok{(}\AttributeTok{x =}\NormalTok{ GDP\_Group, }\AttributeTok{y =}\NormalTok{ Value.x)) }\SpecialCharTok{+}
  \FunctionTok{geom\_boxplot}\NormalTok{() }\SpecialCharTok{+}
  \FunctionTok{labs}\NormalTok{(}
    \AttributeTok{title =} \StringTok{"Sub{-}population Variation: Unemployment in High{-} and Low{-}GDP Countries"}\NormalTok{,}
    \AttributeTok{x =} \StringTok{"GDP per Capita Group"}\NormalTok{,}
    \AttributeTok{y =} \StringTok{"Unemployment Rate (\%)"}
\NormalTok{  ) }\SpecialCharTok{+}
  \FunctionTok{theme\_minimal}\NormalTok{()}
\end{Highlighting}
\end{Shaded}

\pandocbounded{\includegraphics[keepaspectratio]{Template_Assignment_files/Template_Assignment_files/figure-latex/unnamed-chunk-5-1.pdf}}

This boxplot compares the distribution of unemployment rates between
low-GDP and high-GDP countries. The countries are grouped based on
whether their GDP per capita is below or above the median GDP per capita
in the dataset.

The main purpose of the graph is to see whether unemployment rates
differ between richer and poorer countries. If the box for low-GDP
countries is higher, it suggests that countries with lower GDP per
capita generally have higher unemployment. If the high-GDP box is lower,
this suggests that wealthier countries may have more stable labour
markets.

The line inside each box shows the median unemployment rate. The box
itself shows the middle 50\% of the observations. The lines extending
from the box show the wider range of values, while individual dots
represent possible outliers.

\subsection{3.6 Event analysis}\label{event-analysis}

\begin{Shaded}
\begin{Highlighting}[]
\FunctionTok{library}\NormalTok{(dplyr)}
\FunctionTok{library}\NormalTok{(ggplot2)}

\NormalTok{event\_data }\OtherTok{\textless{}{-}}\NormalTok{ merged\_data }\SpecialCharTok{\%\textgreater{}\%}
  \FunctionTok{mutate}\NormalTok{(}
    \AttributeTok{Value.x =} \FunctionTok{as.numeric}\NormalTok{(}\FunctionTok{as.character}\NormalTok{(Value.x)),}
    \AttributeTok{PeriodCode =} \FunctionTok{as.numeric}\NormalTok{(}\FunctionTok{as.character}\NormalTok{(PeriodCode))}
\NormalTok{  ) }\SpecialCharTok{\%\textgreater{}\%}
  \FunctionTok{group\_by}\NormalTok{(PeriodCode) }\SpecialCharTok{\%\textgreater{}\%}
  \FunctionTok{summarise}\NormalTok{(}
    \AttributeTok{Avg\_Unemployment =} \FunctionTok{mean}\NormalTok{(Value.x, }\AttributeTok{na.rm =} \ConstantTok{TRUE}\NormalTok{)}
\NormalTok{  ) }\SpecialCharTok{\%\textgreater{}\%}
  \FunctionTok{filter}\NormalTok{(PeriodCode }\SpecialCharTok{\textgreater{}=} \DecValTok{2006} \SpecialCharTok{\&}\NormalTok{ PeriodCode }\SpecialCharTok{\textless{}=} \DecValTok{2012}\NormalTok{)}

\FunctionTok{ggplot}\NormalTok{(event\_data, }\FunctionTok{aes}\NormalTok{(}\AttributeTok{x =}\NormalTok{ PeriodCode, }\AttributeTok{y =}\NormalTok{ Avg\_Unemployment)) }\SpecialCharTok{+}
  \FunctionTok{geom\_line}\NormalTok{(}\AttributeTok{size =} \DecValTok{1}\NormalTok{) }\SpecialCharTok{+}
  \FunctionTok{geom\_point}\NormalTok{(}\AttributeTok{size =} \DecValTok{3}\NormalTok{) }\SpecialCharTok{+}
  \FunctionTok{geom\_vline}\NormalTok{(}\AttributeTok{xintercept =} \DecValTok{2008}\NormalTok{, }\AttributeTok{linetype =} \StringTok{"dashed"}\NormalTok{, }\AttributeTok{color =} \StringTok{"red"}\NormalTok{) }\SpecialCharTok{+}
  \FunctionTok{annotate}\NormalTok{(}
    \StringTok{"text"}\NormalTok{,}
    \AttributeTok{x =} \DecValTok{2008}\NormalTok{,}
    \AttributeTok{y =} \FunctionTok{max}\NormalTok{(event\_data}\SpecialCharTok{$}\NormalTok{Avg\_Unemployment, }\AttributeTok{na.rm =} \ConstantTok{TRUE}\NormalTok{),}
    \AttributeTok{label =} \StringTok{"2008 financial crisis"}\NormalTok{,}
    \AttributeTok{color =} \StringTok{"red"}\NormalTok{,}
    \AttributeTok{vjust =} \SpecialCharTok{{-}}\FloatTok{0.5}
\NormalTok{  ) }\SpecialCharTok{+}
  \FunctionTok{labs}\NormalTok{(}
    \AttributeTok{title =} \StringTok{"Unemployment Around the 2008 Financial Crisis"}\NormalTok{,}
    \AttributeTok{x =} \StringTok{"Year"}\NormalTok{,}
    \AttributeTok{y =} \StringTok{"Average Unemployment Rate (\%)"}
\NormalTok{  ) }\SpecialCharTok{+}
  \FunctionTok{theme\_minimal}\NormalTok{()}
\end{Highlighting}
\end{Shaded}

\pandocbounded{\includegraphics[keepaspectratio]{Template_Assignment_files/Template_Assignment_files/figure-latex/analysis-1.pdf}}

This figure shows the average unemployment rate around the 2008
financial crisis. After 2008, unemployment increased noticeably in the
countries included in the dataset. This suggests that the crisis may
have had a negative impact on employment. The figure is relevant because
higher unemployment can lead to poverty, financial insecurity, and a
lower quality of life.

\section{Part 4 - Discussion}\label{part-4---discussion}

\subsection{4.1 Discuss your findings}\label{discuss-your-findings}

Our analysis showed a generally negative relationship between GDP per
capita and unemployment. Countries with a higher GDP per capita tended
to have lower unemployment rates. We also observed that unemployment
rates differed considerably between countries, with some countries
consistently experiencing much higher unemployment than others. In
addition, the 2008 financial crisis was followed by a noticeable
increase in unemployment, suggesting that economic shocks can have a
significant impact on labour markets. When GDP falls, economic activity
decreases, causing firms to produce less and often reduce their
workforce, which can lead to higher level of unemployment.

A limitation of this study is that we only analysed data at the national
level. Regional differences within countries were not included, even
though unemployment and economic development can vary strongly between
regions. Future research could examine these regional differences to
gain a more detailed understanding of the relationship between GDP and
unemployment.

\subsection{5.1 Github repository link}\label{github-repository-link}

Provide the link to your PUBLIC repository here:\\
\url{https://github.com/kdamen-lang/tutorial-groep-5-groepje-3}

\subsection{5.2 Reference list}\label{reference-list}

Blanchard, O., Amighini, A., \& Giavazzi, F.
(2021).\emph{Macroeconomics: a European perspective}(4th ed.)
{[}Macroeconomics{]}. Pearson Education
Limited.\url{https://www.pearson.com/uk}

\emph{Gross domestic product (GDP) per capita - Data Portal - United
Nations Economic Commission for Europe}.
(n.d.-b).\url{https://w3.unece.org/PXWeb/en/Table?IndicatorCode=12}

Okun, A. M. (1962). Potential GNP: Its measurement and significance.
\emph{Proceedings of the Business and Economic Statistics Section of the
American Statistical Association}, 98--104.

\emph{Unemployment rate - Data Portal - United Nations Economic
Commission for Europe}.
(n.d.).\url{https://w3.unece.org/PXWeb/en/Table?IndicatorCode=43}\\
\strut \\

\end{document}
